\begin{table*}[!t]\centering\caption{Sustainability Evidence: Measurement Boundary and Claim Strength}\label{tab:susexpanded}\scriptsize\begin{tabularx}{\textwidth}{p{0.14\textwidth}p{0.18\textwidth}p{0.16\textwidth}p{0.12\textwidth}p{0.12\textwidth}Y}\toprule Study/evidence & System boundary & Energy evidence & Direct measurement? & Carbon/LCA? & Defensible interpretation \\\ \midrule Zhao \& You 2023 & adoption/scenario model & modeled energy/emissions & no & scenario carbon & prospective impact, not measured reduction \\\ De Giovanni 2023 & Industry 5.0 model & analytical sustainability mechanism & no & no & modeled trade-offs \\\ Ng et al. 2024 & semantic edge market & energy/cost in real-data simulation & no & no & workload-linked simulation \\\ Chen 2025 & energy-sector review & reported efficiency/consumption evidence & mixed secondary & secondary & review evidence, not direct measurement itself \\\ Neuhaus et al. 2025 & broad sustainability review & predominantly theoretical corpus & no aggregate direct claim & secondary & supports evidence-maturity concern \\\ Bonavolonta et al. 2026 & immersive laboratory & life-cycle assessment & study-specific & yes, LCA & stronger environmental boundary evidence \\\ Mosavisadr et al. 2026 & DT synchronization & sensor-request-count energy proxy & no direct joule claim & no & proxy must not be relabeled as measured energy \\\ \bottomrule\end{tabularx}\end{table*}